\subsection{Интегрирование тригонометрических выражений вида $\int R(\sin x, \cos x)\,dx$. Универсальная тригонометрическая подстановка}

\subsubsection*{1. Постановка задачи}

Рассматривается интеграл
\[
\int R(\sin x,\cos x)\,dx,
\]
где $R(u,v)$ — рациональная функция двух переменных.

\subsubsection*{2. Универсальная тригонометрическая подстановка}

\begin{theorem}[Универсальная подстановка]
Замена $t = \tan\!\dfrac{x}{2}$ рационализирует любой интеграл вида $\int R(\sin x, \cos x)\,dx$.
\end{theorem}

При $t = \tan\dfrac{x}{2}$ выполняются формулы:
\[
\sin x = \frac{2t}{1+t^2}, \qquad
\cos x = \frac{1-t^2}{1+t^2}, \qquad
dx = \frac{2\,dt}{1+t^2}.
\]

После подстановки интеграл принимает вид рационального по $t$:
\[
\int R\!\left(\frac{2t}{1+t^2},\;\frac{1-t^2}{1+t^2}\right)\frac{2\,dt}{1+t^2}.
\]

\textbf{Пример:}
\[
\int\frac{dx}{1+\cos x}
= \int\frac{1}{1+\dfrac{1-t^2}{1+t^2}}\cdot\frac{2\,dt}{1+t^2}
= \int\frac{1}{\dfrac{2}{1+t^2}}\cdot\frac{2\,dt}{1+t^2}
= \int dt = t + C = \tan\frac{x}{2}+C.
\]

\subsubsection*{3. Специальные случаи (более эффективные замены)}

Универсальная подстановка всегда работает, но иногда более эффективны специальные методы.

\textbf{Случай 1.} Если $R(-\sin x, \cos x) = -R(\sin x, \cos x)$ (нечётна по $\sin x$), то:
\[
t = \cos x, \quad dt = -\sin x\,dx.
\]
\textbf{Пример:} $\int \sin^3 x\,dx = \int(1-\cos^2 x)\sin x\,dx \xrightarrow{t=\cos x} -\int(1-t^2)\,dt$.

\textbf{Случай 2.} Если $R(\sin x, -\cos x) = -R(\sin x, \cos x)$ (нечётна по $\cos x$), то:
\[
t = \sin x, \quad dt = \cos x\,dx.
\]
\textbf{Пример:} $\int \cos^3 x\,dx \xrightarrow{t=\sin x} \int(1-t^2)\,dt$.

\textbf{Случай 3.} Если $R(-\sin x, -\cos x) = R(\sin x, \cos x)$ (чётна по обоим аргументам), то:
\[
t = \tan x, \quad dx = \frac{dt}{1+t^2}, \quad
\sin^2 x = \frac{t^2}{1+t^2}, \quad \cos^2 x = \frac{1}{1+t^2}.
\]
\textbf{Пример:} $\int\dfrac{dx}{\sin^2 x \cos^2 x}$.

\subsubsection*{4. Интегрирование $\int \sin^m x\,\cos^n x\,dx$}

\begin{enumerate}
    \item Если $m$ нечётно: выделяем $\sin x\,dx = -d(\cos x)$, ставим $t = \cos x$.
    \item Если $n$ нечётно: выделяем $\cos x\,dx = d(\sin x)$, ставим $t = \sin x$.
    \item Если $m$ и $n$ чётны: используем формулы понижения степени:
    \[
    \sin^2 x = \frac{1-\cos 2x}{2}, \qquad \cos^2 x = \frac{1+\cos 2x}{2},
    \qquad \sin x\cos x = \frac{\sin 2x}{2}.
    \]
\end{enumerate}

\subsubsection*{5. Интегрирование произведений $\sin(mx)\cos(nx)$ и т.\,п.}

Используют формулы произведения:
\[
\sin\alpha\cos\beta = \tfrac{1}{2}[\sin(\alpha-\beta)+\sin(\alpha+\beta)],
\]
\[
\sin\alpha\sin\beta = \tfrac{1}{2}[\cos(\alpha-\beta)-\cos(\alpha+\beta)],
\]
\[
\cos\alpha\cos\beta = \tfrac{1}{2}[\cos(\alpha-\beta)+\cos(\alpha+\beta)].
\]

\subsubsection*{6. Сводная таблица выбора метода}

\begin{center}
\renewcommand{\arraystretch}{1.8}
\begin{tabular}{|l|l|}
\hline
\textbf{Условие} & \textbf{Метод / замена} \\
\hline
Общий случай & $t = \tan(x/2)$ \\
$R$ нечётна по $\sin x$ & $t = \cos x$ \\
$R$ нечётна по $\cos x$ & $t = \sin x$ \\
$R$ чётна по обоим & $t = \tan x$ \\
$\sin^m x\cos^n x$, $m$ или $n$ нечётно & вынести нечётный сомножитель \\
$\sin^m x\cos^n x$, оба чётны & формулы $\cos 2x$ \\
\hline
\end{tabular}
\end{center}